% PAGES: 17-23, 24-30

\noindent Column vectors and row vectors. Column form and row form of a matrix.
\medskip

\noindent Column-based and row-based matrix--vector and matrix--matrix multiplication.
\medskip

\noindent Covariance matrix computation from time series data.
\medskip

\noindent Matrix rank. Nullspace and range of a matrix. Linear independence.
\medskip

\noindent A one period market model example.
\medskip

\noindent Nonsingular matrices and the inverse of a matrix.
\medskip

\noindent Diagonal matrices. Matrix multiplication by diagonal matrices.
\medskip

\noindent Converting between covariance and correlation matrices.
\medskip

\noindent Correlation matrix computation from time series data.
\medskip

\noindent Lower and upper triangular matrices. Tridiagonal matrices.
\medskip

\section{Column and row vectors. Column form and row form of a matrix.}

An $n$-dimensional vector $v \in \R^n$ is denoted by $v=(v_i)_{i=1:n}$ and has $n$
components $v_i \in \R$, for $i=1:n$.\footnote{The vectors and matrices considered
here have entries which are real numbers. While complex numbers will occur
naturally (for example, eigenvalues of a matrices with real entries may be
complex numbers), the presentation and the notations in this book will be
specific to vectors and matrices with real entries.}

The vector $v=(v_i)_{i=1:n}$ is a \textbf{column vector} of size $n$ if

\begin{equation}
v \;=\; \begin{pmatrix} v_1 \\ v_2 \\ \vdots \\ v_n \end{pmatrix}.
\end{equation}

A column vector is also called an $n\times 1$ vector.

The vector $v^t$ is a \textbf{row vector}\footnote{The notation $v^t$ emphasizes
the fact that a row vector is the transpose of a column vector; see also
Definition 1.1.} of size $n$ if

\begin{equation}
v^t \;=\; (v_1\ v_2\ \ldots\ v_n).
\end{equation}

A row vector is also called an $1\times n$ vector.

Unless otherwise specified, a vector $v$ denoted by $v=(v_i)_{i=1:n}$ is a
column vector.

An $m\times n$ matrix $A=(A(j,k))_{j=1:m,k=1:n}$ has $m$ rows and $n$ columns.
Rather than using the entry by entry notation above for the matrix $A$, we
will use either a column-based notation (more often), or a row-based
notation, both being better suited for numerical computations.

The \textbf{column form} of the matrix $A$ is

\begin{equation}
A \;=\; (a_1\ |\ a_2\ |\ \ldots\ |\ a_n) \;=\; \mathrm{col}\,(a_k)_{k=1:n}\,,
\end{equation}

where $a_k$ is the $k$-th column\footnote{For every $k=1:n$, the column
vector $a_k$ is given by $a_k=(A(j,k))_{j=1:m}$.} of $A$, $k=1:n$.

The \textbf{row form} of the matrix $A$ is

\begin{equation}
A \;=\; \begin{pmatrix} r_1 \\ \text{--} \\ r_2 \\ \text{--} \\ \vdots \\ \text{--} \\ r_m \end{pmatrix}
\;=\; \mathrm{row}\,(r_j)_{j=1:m}\,,
\end{equation}

where $r_j$ is the $j$-th row\footnote{For every $j=1:m$, the row vector
$r_j$ is given by $r_j=(A(j,k))_{k=1:n}$.} of $A$, $j=1:m$.

\medskip\noindent\textbf{Row Vector -- Column Vector multiplication:}\footnote{Formula
(1.5) is the same as formula (5.2) for the Euclidean inner product of two
vectors.}\par
Let $v=(v_i)_{i=1:n}$ be a column vector of size $n$, and let
$w^t=(w_1\ w_2\ \ldots\ w_n)$ be a row vector of size $n$. Then,

\begin{equation}
w^t v \;=\; \sum_{i=1}^n w_i v_i.
\end{equation}

\medskip\noindent\textbf{Column Vector -- Row Vector multiplication:}\par
Let $v=(v_j)_{j=1:m}$ be a column vector of size $m$, and let
$w^t=(w_1\ w_2\ \ldots\ w_n)$ be a row vector of size $n$. Then, $vw^t$ is an
$m\times n$ matrix with the following entries:

\begin{equation}
(vw^t)(j,k) \;=\; v_j w_k, \quad \forall\, j=1:m,\ k=1:n.
\end{equation}

\medskip\noindent\textbf{Matrix -- Column Vector multiplication:}\par
Let $A=\mathrm{col}\,(a_k)_{k=1:n}$ be an $m\times n$ matrix given by the
column form (1.3), and let $v=(v_k)_{k=1:n}$ be a column vector of size $n$
given by (1.1). Then,

\begin{equation}
Av \;=\; \sum_{k=1}^n v_k a_k.
\end{equation}

In other words, the result of the multiplication of the column vector $v$ by
the matrix $A$ is a column vector $Av$ which is the linear
combination\footnote{A linear combination of the vectors
$w_1,w_2,\ldots,w_n$ is any sum of these vectors multiplied by real
coefficients, i.e., $c_1w_1+c_2w_2+\ldots+c_nw_n$, where $c_i\in\R$,
$i=1:n$; see also Definition 1.5.} of the columns of $A$ with coefficients
equal to the corresponding entries of $v$.

If $A=\mathrm{row}\,(r_j)_{j=1:m}$ is the row form of $A$, then the $j$-th
entry of the $m\times 1$ column vector $Av$ is

\begin{equation}
(Av)(j) \;=\; r_j v, \quad \forall\, 1\le j \le m.
\end{equation}

Note that, since $r_j$ is a $1\times n$ row vector and $v$ is a $n\times 1$
column vector, it follows from (1.5) that the multiplication from (1.8) can
be performed.

\medskip\noindent\textbf{Row Vector -- Matrix multiplication:}\par
Let $A=\mathrm{row}\,(r_j)_{j=1:m}$ be an $m\times n$ matrix given by the
row form (1.4), and let $w^t=(w_1\ w_2\ \ldots\ w_m)$ be a row vector of
size $m$. Then,

\begin{equation}
w^t A \;=\; \sum_{j=1}^m w_j r_j.
\end{equation}

In other words, the result of the multiplication of the row vector $w^t$ by
the matrix $A$ (from the right) is a row vector $w^tA$ which is the linear
combination of the rows of $A$ with coefficients equal to the corresponding
entries of $w^t$.

If $A=\mathrm{col}\,(a_k)_{k=1:n}$ is the column form of $A$, then the
$k$-th entry of the $1\times n$ row vector $w^tA$ is

\begin{equation}
(w^tA)(k) \;=\; w^t a_k, \quad \forall\, 1\le k\le n.
\end{equation}

\medskip\noindent\textbf{Matrix -- Matrix multiplication:}\par
(i) Let $A$ be an $m\times n$ matrix, and let $B$ be an $n\times p$ matrix
given by $B=\mathrm{col}\,(b_k)_{k=1:p}$. Then, $AB$ is the $m\times p$
matrix given by

\begin{equation}
AB \;=\; \mathrm{col}\,(Ab_k)_{k=1:p} \;=\; (Ab_1\ |\ Ab_2\ |\ \ldots\ |\ Ab_p).
\end{equation}

The result of multiplying the matrices $A$ and $B$ is a matrix whose
columns are the columns of $B$ multiplied by the matrix $A$.

(ii) Let $A$ be an $m\times n$ matrix given by $A=\mathrm{row}\,(r_j)_{j=1:m}$,
and let $B$ be an $n\times p$ matrix. Then, $AB$ is the $m\times p$ matrix
given by

\begin{equation}
AB \;=\; \mathrm{row}\,(r_jB)_{j=1:m} \;=\;
\begin{pmatrix} r_1B \\ \text{--} \\ r_2B \\ \text{--} \\ \vdots \\ \text{--} \\ r_mB \end{pmatrix}.
\end{equation}

The result of multiplying the matrices $A$ and $B$ is a matrix whose rows
are the rows of $A$ multiplied by the matrix $B$.

(iii) Let $A$ be an $m\times n$ matrix given by $A=\mathrm{row}\,(r_j)_{j=1:m}$,
and let $B$ be a $n\times p$ matrix given by $B=\mathrm{col}\,(b_k)_{k=1:p}$.
Then, $AB$ is the $m\times p$ matrix whose entries are given
by\footnote{Note that the multiplication from (1.13) can be performed since
$r_j$ is a $1\times n$ row vector and $b_k$ is a $n\times 1$ column vector;
see (1.5).}

\begin{equation}
(AB)(j,k) \;=\; r_j b_k, \quad \forall\, j=1:m,\ k=1:p.
\end{equation}

\medskip\noindent\textbf{Matrix -- Matrix -- Matrix multiplication:}\par
Let $A$ be an $m\times n$ matrix given by $A=\mathrm{row}\,(r_j)_{j=1:m}$,
let $B$ be an $n\times p$ matrix, and let $C$ be a $p\times l$ matrix given
by $C=\mathrm{col}\,(c_k)_{k=1:l}$. Then, $ABC$ is the $m\times l$ matrix
whose entries are given by

\begin{equation}
(ABC)(j,k) \;=\; r_j B c_k, \quad \forall\, j=1:m,\ k=1:l.
\end{equation}

Note that (1.14) follows from (1.13), since $BC=\mathrm{col}\,(Bc_k)_{k=1:l}$;
cf. (1.11).

Note that matrix multiplication is associative, i.e., $ABC=(AB)C=A(BC)$.

We emphasize again that we almost exclusively think of a matrix as either a
collection of column vectors, or as a collection of row vectors, rather
than as a collection of individual entries. For numerical purposes, this is
an efficient way to implement matrices. Also, linear algebra proofs using
the column form or the row form of a matrix are more insightful and more
compact than proofs using individual entries of a matrix. Most of the
proofs from this book use a vector-based approach.

\begin{definition}
The transpose of an $n\times 1$ column vector $v=(v_i)_{i=1:n}$ is the
$1\times n$ row vector $v^t=(v_1\ v_2\ \ldots\ v_n)$. The transpose of an
$1\times n$ row vector $r=(r_1\ r_2\ \ldots\ r_n)$ is the $n\times 1$ column
vector $r^t=(r_i)_{i=1:n}$.
\end{definition}

Note that

\begin{equation}
(cv)^t \;=\; cv^t, \quad \forall\, v\in\R^n,\ c\in\R.
\end{equation}

\begin{definition}
The transpose matrix $A^t$ of an $m\times n$ matrix $A$ is an $n\times m$
matrix given by
\end{definition}

\begin{equation}
A^t(k,j) \;=\; A(j,k), \quad \forall\, k=1:n,\ j=1:m.
\end{equation}

Transposing a matrix switches the column form of the matrix to a row form,
and the row form of the matrix to a column form as follows:

\begin{align}
A=\mathrm{col}\,(a_k)_{k=1:n} &\iff A^t=\mathrm{row}\,(a_k^t)_{k=1:n}; \\
A=\mathrm{row}\,(r_j)_{j=1:m} &\iff A^t=\mathrm{col}\,(r_j^t)_{j=1:m}.
\end{align}

From (1.16), we find that, for any matrix $A$,

\begin{equation}
(A^t)^t \;=\; A,
\end{equation}

and, for any matrices $A$ and $B$ of the same size,

\begin{equation}
(A+B)^t \;=\; A^t+B^t.
\end{equation}

\begin{lemma}
Let $A$ be an $m\times n$ matrix and let $v$ be a column vector of size
$n$. Then,
\begin{equation}
(Av)^t \;=\; v^tA^t.
\end{equation}
\end{lemma}

\begin{proof}
Let $A=\mathrm{col}\,(a_k)_{k=1:n}$ and $v=(v_i)_{i=1:n}$. Then,
$Av=\sum_{i=1}^n v_ia_i$, and

\begin{equation}
(Av)^t \;=\; \left(\sum_{i=1}^n v_ia_i\right)^t \;=\; \sum_{i=1}^n (v_ia_i)^t
\;=\; \sum_{i=1}^n v_ia_i^t,
\end{equation}

since $v_i\in\R$; see (1.15).

Note that $A^t=\mathrm{row}\,(a_k^t)_{k=1:n}$; cf. (1.17). Then, from (1.9),
it follows that

\begin{equation}
v^tA^t \;=\; \sum_{i=1}^n v_ia_i^t.
\end{equation}

From (1.22) and (1.23), we conclude that $(Av)^t = v^tA^t$.
\end{proof}

It is very important to note that the transpose of the product of two
matrices is \textit{not} the product of the transposes of the two
matrices,\footnote{A similar property holds for inverses of matrices, i.e.,
$(AB)^{-1}\ne A^{-1}B^{-1}$. Moreover, $(AB)^{-1}=B^{-1}A^{-1}$; see Lemma
1.7 for details.} i.e., $(AB)^t\ne A^tB^t$. Instead, the following result
holds:\footnote{The result of Lemma 1.2 extends as follows:
$\left(\prod_{i=1}^{p} A_i\right)^t = \prod_{i=1}^{p} A_{p+1-i}^t$. A proof
can be given by induction; see an exercise at the end of this chapter.}

\begin{lemma}
Let $A$ be an $m\times n$ matrix and let $B$ be an $n\times p$ matrix.
Then,
\begin{equation}
(AB)^t \;=\; B^tA^t.
\end{equation}
\end{lemma}

\begin{proof}
Recall from (1.11) that, if $B=\mathrm{col}\,(b_k)_{k=1:p}$, then
$AB=\mathrm{col}\,(Ab_k)_{k=1:p}$. Thus, from (1.17), we obtain that

\[
(AB)^t \;=\; \left(\mathrm{col}\,(Ab_k)_{k=1:p}\right)^t \;=\;
\mathrm{row}\,\left((Ab_k)^t\right)_{k=1:p}.
\]

Using (1.21), (1.12), and the fact that $B^t=\mathrm{row}\,(b_k^t)_{k=1:p}$,
see (1.17) we conclude that

\begin{equation*}
(AB)^t \;=\; \mathrm{row}\,\left(b_k^tA^t\right)_{k=1:p} \;=\;
\left(\mathrm{row}\,(b_k^t)_{k=1:p}\right)A^t \;=\; B^tA^t. \qedhere
\end{equation*}
\end{proof}

\begin{definition}
A matrix with the same number of rows and columns is called a
\textbf{square matrix}.
\end{definition}

Note that an $n\times n$ square matrix is also called a square matrix of
size $n$.

\begin{definition}
A square matrix is \textbf{symmetric} if and only if the matrix and its
transpose are the same. In other words, a square matrix $A$ of size $n$ is
symmetric if and only if $A=A^t$, i.e.,
\[
A(j,k) \;=\; A(k,j), \quad \forall\, 1\le j<k\le n;
\]
\end{definition}

The product of two symmetric matrices is not necessarily a symmetric
matrix, as seen in the example below.

\begin{examplex}
Let $A=\begin{pmatrix}1&2\\2&1\end{pmatrix}$ and
$B=\begin{pmatrix}2&1\\1&0\end{pmatrix}$ be two symmetric matrices. Then,
\[
AB \;=\; \begin{pmatrix}4&1\\5&2\end{pmatrix} \;\ne\; (AB)^t \;=\;
\begin{pmatrix}4&5\\1&2\end{pmatrix}
\]
\end{examplex}

The identity matrix,\footnote{The $n\times n$ identity matrix is sometimes
denoted by $I_n$. We do not use this notation, but denote by $I$ identity
matrices of any size.} denoted by $I$, is a square matrix with entries
equal to 1 on the main diagonal and equal to 0 everywhere else, i.e.,

\[
I \;=\; \begin{pmatrix} 1 & 0 & \ldots & 0 \\ 0 & 1 & \ldots & 0 \\
\vdots & \vdots & \ddots & \vdots \\ 0 & 0 & \ldots & 1 \end{pmatrix}.
\]

The $k$-th column of the identity matrix is denoted by $e_k$. Thus,

\begin{equation}
e_k(i) = 0, \text{ for } 1\le i\ne k\le n \text{ and } e_k(k)=1.
\end{equation}

The column form and row form of the identity matrix $I$ are, respectively,

\[
I \;=\; \mathrm{col}\,(e_k)_{k=1:n}; \quad I \;=\; \mathrm{row}\,(e_k^t)_{k=1:n};
\]

cf. (1.17), since $I=I^t$.

\begin{lemma}
(i) Let $A=\mathrm{col}\,(a_k)_{k=1:n}$ be an $m\times n$ matrix. If $e_k$
is the $k$-th column of the $n\times n$ identity matrix, then

\begin{equation}
Ae_k \;=\; a_k, \quad \forall\, k=1:n,
\end{equation}

and therefore $AI=A$.

(ii) Let $A=\mathrm{row}\,(r_j)_{j=1:m}$ be an $m\times n$ matrix. If $e_j$
is the $j$-th column of the $m\times m$ identity matrix, then

\begin{equation}
e_j^tA \;=\; r_j, \quad \forall\, j=1:m,
\end{equation}

and therefore $IA=A$.
\end{lemma}

\begin{proof}
(i) Let $A=\mathrm{col}\,(a_k)_{k=1:n}$. Recall from (1.25) that
$e_k(k)=1$ and $e_k(i)=0$, for $i\ne k$. From (1.7), we obtain that

\begin{equation}
Ae_k \;=\; \sum_{i=1}^n e_k(i)a_i \;=\; a_k.
\end{equation}

If $I=\mathrm{col}\,(e_k)_{k=1:n}$, it follows from (1.11) and (1.28) that

\[
AI \;=\; \mathrm{col}\,(Ae_k)_{k=1:n} \;=\; \mathrm{col}\,(a_k)_{k=1:n} \;=\; A.
\]

(ii) Let $A=\mathrm{row}\,(r_j)_{j=1:m}$. Recall from (1.25) that
$e_j(j)=1$ and $e_j(i)=0$, for $i\ne j$. From (1.9), we find that

\begin{equation}
e_j^tA \;=\; \sum_{i=1}^m e_j(i)r_i \;=\; r_j.
\end{equation}

If $I=\mathrm{row}\,(e_j^t)_{j=1:m}$, it follows from (1.12) and (1.29)
that

\begin{equation*}
IA \;=\; \mathrm{row}\,(e_j^tA)_{j=1:m} \;=\; \mathrm{row}\,(r_j)_{j=1:m} \;=\; A. \qedhere
\end{equation*}
\end{proof}

\subsection{Covariance matrix computation from time series data}

Let $X_1, X_2, \ldots, X_n$ be random variables given by time series data
at $N$ data points $t_i$, $i=1:N$. In other words, the values of $X_k(t_i)$
are given for all $k=1:n$ and $i=1:N$.

Denote by $\widehat{\mu}_{X_k}$ the sample mean of the random variable
$X_k$, for $k=1:n$, i.e.,

\[
\widehat{\mu}_{X_k} \;=\; \frac{1}{N}\sum_{i=1}^N X_k(t_i).
\]

The sample covariance matrix $\widehat{\Sigma}_{\bfX}$ of the $n$ random
variables $X_1, X_2, \ldots, X_n$ is the $n\times n$ square matrix with
entries

\begin{equation}
\widehat{\Sigma}_{\bfX}(j,k) \;=\; \widehat{\cov}(X_j,X_k), \quad \forall\, 1\le j,k\le n,
\end{equation}

where $\widehat{\cov}(X_j,X_k)$ is the unbiased sample covariance of the
random variables $X_j$ and $X_k$ given by

\begin{equation}
\widehat{\cov}(X_j,X_k) \;=\; \frac{1}{N-1}\sum_{i=1}^N
(X_j(t_i)-\widehat{\mu}_{X_j})(X_k(t_i)-\widehat{\mu}_{X_k}).
\end{equation}

From (1.30) and (1.31), we find that

\begin{equation}
\widehat{\Sigma}_{\bfX}(j,k) \;=\; \frac{1}{N-1}\sum_{i=1}^N
(X_j(t_i)-\widehat{\mu}_{X_j})(X_k(t_i)-\widehat{\mu}_{X_k}).
\end{equation}

The sample covariance matrix $\widehat{\Sigma}_{\bfX}$ is symmetric since,
from (1.32), it follows that

\begin{align*}
\widehat{\Sigma}_{\bfX}(j,k)
&= \frac{1}{N-1}\sum_{i=1}^N (X_j(t_i)-\widehat{\mu}_{X_j})(X_k(t_i)-\widehat{\mu}_{X_k}) \\
&= \frac{1}{N-1}\sum_{i=1}^N (X_k(t_i)-\widehat{\mu}_{X_k})(X_j(t_i)-\widehat{\mu}_{X_j}) \\
&= \widehat{\Sigma}_{\bfX}(k,j), \quad \forall\, 1\le j,k\le n.
\end{align*}

The sample covariance matrix can be computed efficiently by using matrix
formulation for the time series data $X_k(t_i)$, $i=1:N$, $k=1:n$, as shown
below.

Let $T_{\bfX}$ be the corresponding $N\times n$ matrix of time series
data, i.e., let $T_{\bfX}=(T_{\bfX}(i,k))_{i=1:N,k=1:n}$ with

\begin{equation}
T_{\bfX}(i,k) \;=\; X_k(t_i), \quad \forall\, 1\le k\le n,\ 1\le i\le N.
\end{equation}

Let $\overline{T}_{\bfX}$ be the $N\times n$ matrix of time series data where
the sample mean of each random variable is subtracted from the
corresponding time series data, i.e., let $\overline{T}_{\bfX} =
(\overline{T}_{\bfX}(i,k))_{i=1:N,k=1:n}$ with

\begin{equation}
\overline{T}_{\bfX}(i,k) \;=\; X_k(t_i) - \widehat{\mu}_{X_k}, \quad
\forall\, 1\le k\le n,\ 1\le i\le N.
\end{equation}

Then, the sample covariance matrix $\widehat{\Sigma}_{\bfX}$ can be computed
from $\overline{T}_{\bfX}$ as follows:

\begin{equation}
\widehat{\Sigma}_{\bfX} \;=\; \frac{1}{N-1}\,\overline{T}_{\bfX}^t\overline{T}_{\bfX}.
\end{equation}

For clarity, we include an example below and the proof of (1.35).

\medskip
\noindent\rule{\linewidth}{0.4pt}
\medskip

\noindent\textit{Example:}\ The end of day adjusted close prices for Apple,
Facebook, Google, Microsoft, and Yahoo between 1/10/2013 and 1/29/2013 were:

\begin{center}
\begin{tabular}{cccccc}
Date & AAPL & FB & GOOG & MSFT & YHOO \\[2pt]
1/10/2013 & 523.51 & 31.30 & 741.48 & 26.46 & 18.99 \\
1/11/2013 & 520.30 & 31.72 & 739.99 & 26.83 & 19.29 \\
1/14/2013 & 501.75 & 30.95 & 723.25 & 26.89 & 19.43 \\
1/15/2013 & 485.92 & 30.10 & 724.93 & 27.21 & 19.52 \\
1/16/2013 & 506.09 & 29.85 & 715.19 & 27.04 & 20.07 \\
1/17/2013 & 502.68 & 30.14 & 711.32 & 27.25 & 20.13 \\
1/18/2013 & 500.00 & 29.66 & 704.51 & 27.25 & 20.02 \\
1/22/2013 & 504.77 & 30.73 & 702.87 & 27.15 & 19.90 \\
1/23/2013 & 514.01 & 30.82 & 741.50 & 27.61 & 20.11 \\
1/24/2013 & 450.50 & 31.08 & 754.21 & 27.63 & 20.44 \\
1/25/2013 & 439.88 & 31.54 & 753.67 & 27.88 & 20.37 \\
1/28/2013 & 449.83 & 32.47 & 750.73 & 27.91 & 20.31 \\
1/29/2013 & 458.27 & 30.79 & 753.68 & 28.01 & 19.70 \\
\end{tabular}
\end{center}

The time series matrix of the daily returns\footnote{Unless specified
otherwise, the return between times $\tau_1$ and $\tau_2$ of an asset with
spot prices $S(\tau_1)$ and $S(\tau_2)$ will mean the percentage return,
which is $\frac{S(\tau_2)-S(\tau_1)}{S(\tau_1)}$.} of the five stocks above
between 1/11/2013 and 1/29/2013 is

\[
T_{\bfX} \;=\;
\begin{pmatrix}
-0.0061 & 0.0134 & -0.0020 & 0.0140 & 0.0158 \\
-0.0357 & -0.0243 & -0.0226 & 0.0022 & 0.0073 \\
-0.0315 & -0.0275 & 0.0023 & 0.0119 & 0.0046 \\
0.0415 & -0.0083 & -0.0134 & -0.0062 & 0.0282 \\
-0.0067 & 0.0097 & -0.0054 & 0.0078 & 0.0030 \\
-0.0053 & -0.0159 & -0.0096 & 0.0000 & -0.0055 \\
0.0095 & 0.0361 & -0.0023 & -0.0037 & -0.0060 \\
0.0183 & 0.0029 & 0.0550 & 0.0169 & 0.0106 \\
-0.1236 & 0.0084 & 0.0171 & 0.0007 & 0.0164 \\
-0.0236 & 0.0148 & -0.0007 & 0.0090 & -0.0034 \\
0.0226 & 0.0295 & -0.0039 & 0.0011 & -0.0029 \\
0.0188 & -0.0517 & 0.0039 & 0.0036 & -0.0300
\end{pmatrix},
\]

where, e.g., the daily return of GOOG on 1/24/2013 is

\[
\frac{754.21-741.50}{741.50} \;=\; 0.0171 \;=\; T_{\bfX}(9,3),
\]

and the daily return of GOOG on 1/28/2013 is

\[
\frac{750.73-753.67}{753.67} \;=\; -0.0039 \;=\; T_{\bfX}(11,3).
\]

The sample means of the returns of the five stocks are $-0.0101$ (AAPL),
$-0.0011$ (FB), $0.0015$ (GOOG), $0.0048$ (MSFT), and $0.0032$ (YHOO). By
subtracting the sample mean of each column of $T_{\bfX}$ we obtain from
(1.34) that

\begin{equation}
\overline{T}_{\bfX} \;=\;
\begin{pmatrix}
0.0040 & 0.0145 & -0.0035 & 0.0092 & 0.0126 \\
-0.0255 & -0.0232 & -0.0242 & -0.0025 & 0.0041 \\
-0.0214 & -0.0264 & 0.0008 & 0.0071 & 0.0015 \\
0.0517 & -0.0072 & -0.0150 & -0.0110 & 0.0250 \\
0.0034 & 0.0108 & -0.0069 & 0.0030 & -0.0002 \\
0.0048 & -0.0149 & -0.0111 & -0.0048 & -0.0086 \\
0.0197 & 0.0371 & -0.0039 & -0.0084 & -0.0092 \\
0.0285 & 0.0040 & 0.0534 & 0.0122 & 0.0074 \\
-0.1134 & 0.0095 & 0.0156 & -0.0041 & 0.0132 \\
-0.0134 & 0.0159 & -0.0022 & 0.0043 & -0.0066 \\
0.0328 & 0.0306 & -0.0054 & -0.0037 & -0.0061 \\
0.0289 & -0.0507 & 0.0024 & -0.0012 & -0.0332
\end{pmatrix}.
\end{equation}

\medskip
\noindent\rule{\linewidth}{0.4pt}
\medskip

We now show that the formula (1.35) holds, i.e.,

\begin{equation}
\widehat{\Sigma}_{\bfX} \;=\; \frac{1}{N-1}\,\overline{T}_{\bfX}^t\overline{T}_{\bfX};
\end{equation}

see also Theorem 7.1 and the proof therein.

From (1.34), we find that, for any $1\le j,k\le n$,

\begin{equation}
\overline{T}_{\bfX}(i,k) \;=\; X_k(t_i) - \widehat{\mu}_{X_k} \quad
\text{and} \quad \overline{T}_{\bfX}(i,j) \;=\; X_j(t_i) - \widehat{\mu}_{X_j},
\quad \forall\, i=1:N.
\end{equation}

Then, from (1.32) and (1.38), it follows that

\begin{align}
\widehat{\Sigma}_{\bfX}(j,k) &= \frac{1}{N-1}\sum_{i=1}^N
(X_j(t_i)-\widehat{\mu}_{X_j})(X_k(t_i)-\widehat{\mu}_{X_k}) \\
&= \frac{1}{N-1}\sum_{i=1}^N \overline{T}_{\bfX}(i,j)\overline{T}_{\bfX}(i,k),
\quad \forall\, 1\le j,k\le n.
\end{align}

Let $\overline{T}_{X_k}$ be the $N\times 1$ column vector of the time series
data for the random variable $X_k$ with $\widehat{\mu}_{X_k}$ subtracted
from each data value, i.e.,

\[
\overline{T}_{X_k} \;=\; \left(X_k(t_i) - \widehat{\mu}_{X_k}\right)_{i=1:N}.
\]

The time series matrix $\overline{T}_{\bfX} =
(\overline{T}_{\bfX}(i,k))_{i=1:N,k=1:n}$ has the following column form:

\[
\overline{T}_{\bfX} \;=\; \mathrm{col}\,\left(\overline{T}_{X_k}\right)_{k=1:n}.
\]

Moreover,

\[
\overline{T}_{\bfX}(i,j) \;=\; \overline{T}_{X_j}(i) \quad \text{and} \quad
\overline{T}_{\bfX}(i,k) \;=\; \overline{T}_{X_k}(i), \quad \forall\,
1\le j,k\le n,\ 1\le i\le N,
\]

and, from (1.40), we obtain that

\begin{align}
\widehat{\Sigma}_{\bfX}(j,k) &= \frac{1}{N-1}\sum_{i=1}^N
\overline{T}_{X_j}(i)\overline{T}_{X_k}(i) \\
&= \frac{1}{N-1}\,\overline{T}_{X_j}^t\overline{T}_{X_k}, \quad \forall\,
1\le j,k\le n,
\end{align}

where the last equality follows from the row vector--column vector
multiplication formula (1.5).

Since $\overline{T}_{\bfX} = \mathrm{col}\,(\overline{T}_{X_k})_{k=1:n}$, it
follows that $\overline{T}_{\bfX}^t = \mathrm{row}\,(\overline{T}_{X_j}^t)_{j=1:n}$;
see (1.17). From (1.13), we obtain that the $(j,k)$ entry of the matrix
$\overline{T}_{\bfX}^t\overline{T}_{\bfX}$ is

\begin{equation}
(\overline{T}_{\bfX}^t\overline{T}_{\bfX})(j,k) \;=\;
\overline{T}_{X_j}^t\overline{T}_{X_k}, \quad \forall\, 1\le j,k\le n.
\end{equation}

Then, from (1.42) and (1.43), we conclude that

\[
\widehat{\Sigma}_{\bfX}(j,k) \;=\; \frac{1}{N-1}\,
(\overline{T}_{\bfX}^t\overline{T}_{\bfX})(j,k), \quad \forall\, 1\le j,k\le n,
\]

and therefore

\[
\widehat{\Sigma}_{\bfX} \;=\; \frac{1}{N-1}\,\overline{T}_{\bfX}^t\overline{T}_{\bfX},
\]

which is what we wanted to prove; see (1.37).

\medskip
\noindent\rule{\linewidth}{0.4pt}
\medskip

\noindent\textit{Example (continued):}\ The sample covariance matrix
$\widehat{\Sigma}_{\bfX}$ of the daily returns of AAPL, FB, GOOG, MSFT, YHOO
between 1/11/2013 and 1/29/2013 can be computed using formula (1.35) with
$N=12$ and $\overline{T}_{\bfX}$ given by (1.36). We find that

\begin{equation}
\widehat{\Sigma}_{\bfX} \;=\;
\begin{pmatrix}
0.0018 & 0.0000 & -0.0001 & 0.0000 & -0.0001 \\
0.0000 & 0.0006 & 0.0001 & 0.0000 & 0.0001 \\
-0.0001 & 0.0001 & 0.0004 & 0.0001 & 0.0000 \\
0.0000 & 0.0000 & 0.0001 & 0.0001 & 0.0000 \\
-0.0001 & 0.0001 & 0.0000 & 0.0000 & 0.0002
\end{pmatrix}. \qquad \square
\end{equation}

\medskip
\noindent\rule{\linewidth}{0.4pt}
\medskip

More properties of covariance matrices obtained from time series data can
be found in section 7.2.
